%% snippets.tex --- LaTeX-ready blocks for the humanization pass.
%% Self-contained: paste the pieces you want. Uses only standard article +
%% amsmath; the boxes use \fbox/minipage so no extra packages are required.
%% Check by appending the "demo" document at the end and running pdflatex.

% ---------------------------------------------------------------- abstract --
\begin{abstract}
Someone publishes a number in years --- an aquifer with a decade or two left, a
phosphate reserve of three centuries --- and a reader hears a date. Three
different quantities travel under that single label: a ratio of an economic
reserve classification to current production, a trend fitted to a satellite
gravity anomaly, and a log biomass margin divided by a fishing mortality. Each
answers a real question; none is a time at which anything physically runs out.
A weighted aggregate compounds the confusion by hiding a component deficit
behind a positive scalar.

We build the accounting layer that keeps the three apart. Compartments carry a
material identity, a boundary, and a unit, so biomass, money, and biodiversity
are never summed into one scalar; conservation follows from the incidence
structure of the flux network, positivity from donor limitation, and services
are readouts rather than conserved mass. We separate three predicates the
literature conflates --- accounting consistency, stoichiometric conservation,
and barrier admissibility --- and prove their relations, so that each claim about
a material system carries the predicate it actually establishes. Depletion time
splits into gross turnover intensity, a frozen-rate ratio, and a
scenario-conditioned hitting time, with uniform-drift bounds between them. No
nonnegative weighting of component balances certifies that every component
clears its floor: certification needs the vector. We classify three public
applications at their exact status and state the surrogate on which each
first-passage result is conditioned. We name the gaps that remain open, and we
read weak and strong sustainability as two regimes of one system, distinguished
by whether the material cycle closes at the rate of use.

\medskip
\noindent\textbf{Keywords:} material flow accounting; stock--flow ledger;
depletion indicators; first-passage time; conservation laws; composite
indicators; reserve life
\end{abstract}

% --------------------------------------------------------- reader's note --
\noindent\fbox{\begin{minipage}{0.96\textwidth}\small
\textbf{Who this paper is for, and in what order.} The verdicts, in thirty
seconds: the title, this note, and \S\ref{sec:applied}. The argument, in three
minutes: \S1, \S6, \S10.1, \S11. The machinery, in an hour: \S2 (the ledger),
\S3 (three certificates), \S4 (what the closed ledger guarantees), \S5
(readouts and the componentwise deficit), \S6--\S7 (depletion time and its
passage semantics), \S8 (domain templates at registered status), \S9 (the
interface with institutional delay dynamics). Readers who want only the applied
conclusion can take \S6.5 and the classification table; the theorems are what
license those readings, and nothing else in this paper is required to accept
them.

\smallskip
\textbf{What this paper is not.} It is not a forecast, not an index, and not a
recommendation about any particular aquifer, fishery, or mine. It is an
accounting layer whose purpose is to make sure that when a statement about
depletion is made, the statement that is true is the statement that is written
down.
\end{minipage}}

% ------------------------------------------------- opening of \S1.1 (new) --
\subsection{Two ways the bookkeeping goes wrong}\label{sec:twoways}
A water board announces that a district has a decade or two of groundwater at
current rates. A geological survey lists phosphate reserves of roughly
$1{,}000{,}000$\,kt, and someone divides by production and announces three
centuries. A fisheries assessment divides a log biomass margin by a fishing
mortality and calls the quotient a time scale. All three numbers are computed
correctly. All three answer different questions, and only the last is even
described in the language of its own construction.

Sustainability accounting fails in two characteristic ways. The first is
\emph{compensatory aggregation}: heterogeneous physical stocks and service
flows are summarised by scalar indices whose cross-component trades are never
declared as mathematics, so a severe deficit in one component can coexist with
a positive aggregate. The composite-indicator and
weak-versus-strong-sustainability literatures document the failure and its
noncompensatory remedies (Munda and Nardo, 2009; Ekins et al., 2003; Neumayer,
2013). The second is \emph{classification drift}: quantities carrying the units
of time circulate as if they were one quantity, ``time to depletion'', though
they answer different questions under different assumptions.

We do not dispute the arithmetic of any of these indicators. We dispute the
shared label. The two theses are not equivalent, and only the second is
defensible for a ledger of this kind.

% --------------------------------------------------- gaps named early, \S1 --
\paragraph{Four limits, registered rather than conceded later.}
(i) The groundwater two-pool model is an open gap; the applied object we admit
is its one-pool affine approximation (\S8.2). (ii) Spawning biomass is not an
abiotic support pool, and we do not treat it as one (\S6.5.4). (iii) Every
first-passage result is conditioned on a declared surrogate and carries
explicit non-claims (\S7.7). (iv) The numerical exhibits are worked instances of
the constructions, not constitutive claims for the named domains. A reader
looking for guidance about a specific resource will not find it here; a reader
looking to know which of their claims their data can carry will.

% --------------------------------------------------- section prefaces (\S) --
% One italic sentence per section; adjust \ref labels to your own.
\newcommand{\preface}[1]{\emph{#1}\par\smallskip}

\preface{We fix the objects first: what a compartment is, what a primitive flux
is, and why a conversion may appear only as an explicit coefficient. Everything
that follows is a property of this bookkeeping and nothing else.}   % \S2

\preface{Three predicates get conflated in the literature: the ledger balances,
the chemistry closes, the barriers hold. We prove each separately so that a
reader can tell which one a given argument uses.}                   % \S3

\preface{Here the ledger pays for itself. Mass is conserved because of the
incidence pattern, nothing goes negative because every outflow is
donor-limited, and no positive-effort interior rest exists to hide in.}  % \S4

\preface{Services are read off the ledger, not added to it. The deficit is a
vector, and \S5.4 shows exactly where a scalar would have hidden it.}  % \S5

\preface{Three quantities, three questions. We keep them named separately through
\S6.5, where each public indicator is classified at the status its construction
supports.}                                                             % \S6

\preface{A hitting time is a property of a model, so we state the model, the
barrier, and the conditioning, and then say what does not follow (\S7.7).}  % \S7

\preface{Templates for phosphorus, groundwater, and harvest-side economics, each
with its status and its admitted object marked, so a reader can see what is a
construction and what is a claim about the world.}                    % \S8

\preface{One shared object, two systems, and a non-reduction result that says
why the connection is an interface contract rather than a limit.}     % \S9

\preface{Not caveats: prohibitions on a class of argument. Five double-counting
rules, a theorem against compensatory weighting, and a discipline that keeps
negative and boundary content first-class.}                          % \S10

% --------------------------------------- reusable verdict sentences (\S6.5) --
\newcommand{\verdictGThreeP}{A trend index on a gravity anomaly is a statistical
statement about a record, not a stock ratio.}
\newcommand{\verdictPhosphate}{Reserves divided by production is arithmetic on an
economic classification, not a forecast; the classification refills itself as it
is drawn down.}
\newcommand{\verdictFisheries}{A log margin divided by a removal rate is a
pressure scale, not a time to any event.}
% Use \verdictGThreeP\ in \S1, in the \S6.5 heading text, and again in \S11.

% ------------------------------------------------------ \S11 closing lines --
Nothing in this paper argues that the numbers are wrong. It argues that a reader
should be able to tell, from the statement itself, which question was answered.
That is a bookkeeping property, and it is cheap to install.

%% ------------------------------------------------------------------- demo --
%% \documentclass{article}
%% \usepackage{amsmath,amssymb}
%% \newlabel{sec:applied}{{6.5}{0}}
%% \begin{document}
%% %% snippets.tex --- LaTeX-ready blocks for the humanization pass.
%% Self-contained: paste the pieces you want. Uses only standard article +
%% amsmath; the boxes use \fbox/minipage so no extra packages are required.
%% Check by appending the "demo" document at the end and running pdflatex.

% ---------------------------------------------------------------- abstract --
\begin{abstract}
Someone publishes a number in years --- an aquifer with a decade or two left, a
phosphate reserve of three centuries --- and a reader hears a date. Three
different quantities travel under that single label: a ratio of an economic
reserve classification to current production, a trend fitted to a satellite
gravity anomaly, and a log biomass margin divided by a fishing mortality. Each
answers a real question; none is a time at which anything physically runs out.
A weighted aggregate compounds the confusion by hiding a component deficit
behind a positive scalar.

We build the accounting layer that keeps the three apart. Compartments carry a
material identity, a boundary, and a unit, so biomass, money, and biodiversity
are never summed into one scalar; conservation follows from the incidence
structure of the flux network, positivity from donor limitation, and services
are readouts rather than conserved mass. We separate three predicates the
literature conflates --- accounting consistency, stoichiometric conservation,
and barrier admissibility --- and prove their relations, so that each claim about
a material system carries the predicate it actually establishes. Depletion time
splits into gross turnover intensity, a frozen-rate ratio, and a
scenario-conditioned hitting time, with uniform-drift bounds between them. No
nonnegative weighting of component balances certifies that every component
clears its floor: certification needs the vector. We classify three public
applications at their exact status and state the surrogate on which each
first-passage result is conditioned. We name the gaps that remain open, and we
read weak and strong sustainability as two regimes of one system, distinguished
by whether the material cycle closes at the rate of use.

\medskip
\noindent\textbf{Keywords:} material flow accounting; stock--flow ledger;
depletion indicators; first-passage time; conservation laws; composite
indicators; reserve life
\end{abstract}

% --------------------------------------------------------- reader's note --
\noindent\fbox{\begin{minipage}{0.96\textwidth}\small
\textbf{Who this paper is for, and in what order.} The verdicts, in thirty
seconds: the title, this note, and \S\ref{sec:applied}. The argument, in three
minutes: \S1, \S6, \S10.1, \S11. The machinery, in an hour: \S2 (the ledger),
\S3 (three certificates), \S4 (what the closed ledger guarantees), \S5
(readouts and the componentwise deficit), \S6--\S7 (depletion time and its
passage semantics), \S8 (domain templates at registered status), \S9 (the
interface with institutional delay dynamics). Readers who want only the applied
conclusion can take \S6.5 and the classification table; the theorems are what
license those readings, and nothing else in this paper is required to accept
them.

\smallskip
\textbf{What this paper is not.} It is not a forecast, not an index, and not a
recommendation about any particular aquifer, fishery, or mine. It is an
accounting layer whose purpose is to make sure that when a statement about
depletion is made, the statement that is true is the statement that is written
down.
\end{minipage}}

% ------------------------------------------------- opening of \S1.1 (new) --
\subsection{Two ways the bookkeeping goes wrong}\label{sec:twoways}
A water board announces that a district has a decade or two of groundwater at
current rates. A geological survey lists phosphate reserves of roughly
$1{,}000{,}000$\,kt, and someone divides by production and announces three
centuries. A fisheries assessment divides a log biomass margin by a fishing
mortality and calls the quotient a time scale. All three numbers are computed
correctly. All three answer different questions, and only the last is even
described in the language of its own construction.

Sustainability accounting fails in two characteristic ways. The first is
\emph{compensatory aggregation}: heterogeneous physical stocks and service
flows are summarised by scalar indices whose cross-component trades are never
declared as mathematics, so a severe deficit in one component can coexist with
a positive aggregate. The composite-indicator and
weak-versus-strong-sustainability literatures document the failure and its
noncompensatory remedies (Munda and Nardo, 2009; Ekins et al., 2003; Neumayer,
2013). The second is \emph{classification drift}: quantities carrying the units
of time circulate as if they were one quantity, ``time to depletion'', though
they answer different questions under different assumptions.

We do not dispute the arithmetic of any of these indicators. We dispute the
shared label. The two theses are not equivalent, and only the second is
defensible for a ledger of this kind.

% --------------------------------------------------- gaps named early, \S1 --
\paragraph{Four limits, registered rather than conceded later.}
(i) The groundwater two-pool model is an open gap; the applied object we admit
is its one-pool affine approximation (\S8.2). (ii) Spawning biomass is not an
abiotic support pool, and we do not treat it as one (\S6.5.4). (iii) Every
first-passage result is conditioned on a declared surrogate and carries
explicit non-claims (\S7.7). (iv) The numerical exhibits are worked instances of
the constructions, not constitutive claims for the named domains. A reader
looking for guidance about a specific resource will not find it here; a reader
looking to know which of their claims their data can carry will.

% --------------------------------------------------- section prefaces (\S) --
% One italic sentence per section; adjust \ref labels to your own.
\newcommand{\preface}[1]{\emph{#1}\par\smallskip}

\preface{We fix the objects first: what a compartment is, what a primitive flux
is, and why a conversion may appear only as an explicit coefficient. Everything
that follows is a property of this bookkeeping and nothing else.}   % \S2

\preface{Three predicates get conflated in the literature: the ledger balances,
the chemistry closes, the barriers hold. We prove each separately so that a
reader can tell which one a given argument uses.}                   % \S3

\preface{Here the ledger pays for itself. Mass is conserved because of the
incidence pattern, nothing goes negative because every outflow is
donor-limited, and no positive-effort interior rest exists to hide in.}  % \S4

\preface{Services are read off the ledger, not added to it. The deficit is a
vector, and \S5.4 shows exactly where a scalar would have hidden it.}  % \S5

\preface{Three quantities, three questions. We keep them named separately through
\S6.5, where each public indicator is classified at the status its construction
supports.}                                                             % \S6

\preface{A hitting time is a property of a model, so we state the model, the
barrier, and the conditioning, and then say what does not follow (\S7.7).}  % \S7

\preface{Templates for phosphorus, groundwater, and harvest-side economics, each
with its status and its admitted object marked, so a reader can see what is a
construction and what is a claim about the world.}                    % \S8

\preface{One shared object, two systems, and a non-reduction result that says
why the connection is an interface contract rather than a limit.}     % \S9

\preface{Not caveats: prohibitions on a class of argument. Five double-counting
rules, a theorem against compensatory weighting, and a discipline that keeps
negative and boundary content first-class.}                          % \S10

% --------------------------------------- reusable verdict sentences (\S6.5) --
\newcommand{\verdictGThreeP}{A trend index on a gravity anomaly is a statistical
statement about a record, not a stock ratio.}
\newcommand{\verdictPhosphate}{Reserves divided by production is arithmetic on an
economic classification, not a forecast; the classification refills itself as it
is drawn down.}
\newcommand{\verdictFisheries}{A log margin divided by a removal rate is a
pressure scale, not a time to any event.}
% Use \verdictGThreeP\ in \S1, in the \S6.5 heading text, and again in \S11.

% ------------------------------------------------------ \S11 closing lines --
Nothing in this paper argues that the numbers are wrong. It argues that a reader
should be able to tell, from the statement itself, which question was answered.
That is a bookkeeping property, and it is cheap to install.

%% ------------------------------------------------------------------- demo --
%% \documentclass{article}
%% \usepackage{amsmath,amssymb}
%% \newlabel{sec:applied}{{6.5}{0}}
%% \begin{document}
%% %% snippets.tex --- LaTeX-ready blocks for the humanization pass.
%% Self-contained: paste the pieces you want. Uses only standard article +
%% amsmath; the boxes use \fbox/minipage so no extra packages are required.
%% Check by appending the "demo" document at the end and running pdflatex.

% ---------------------------------------------------------------- abstract --
\begin{abstract}
Someone publishes a number in years --- an aquifer with a decade or two left, a
phosphate reserve of three centuries --- and a reader hears a date. Three
different quantities travel under that single label: a ratio of an economic
reserve classification to current production, a trend fitted to a satellite
gravity anomaly, and a log biomass margin divided by a fishing mortality. Each
answers a real question; none is a time at which anything physically runs out.
A weighted aggregate compounds the confusion by hiding a component deficit
behind a positive scalar.

We build the accounting layer that keeps the three apart. Compartments carry a
material identity, a boundary, and a unit, so biomass, money, and biodiversity
are never summed into one scalar; conservation follows from the incidence
structure of the flux network, positivity from donor limitation, and services
are readouts rather than conserved mass. We separate three predicates the
literature conflates --- accounting consistency, stoichiometric conservation,
and barrier admissibility --- and prove their relations, so that each claim about
a material system carries the predicate it actually establishes. Depletion time
splits into gross turnover intensity, a frozen-rate ratio, and a
scenario-conditioned hitting time, with uniform-drift bounds between them. No
nonnegative weighting of component balances certifies that every component
clears its floor: certification needs the vector. We classify three public
applications at their exact status and state the surrogate on which each
first-passage result is conditioned. We name the gaps that remain open, and we
read weak and strong sustainability as two regimes of one system, distinguished
by whether the material cycle closes at the rate of use.

\medskip
\noindent\textbf{Keywords:} material flow accounting; stock--flow ledger;
depletion indicators; first-passage time; conservation laws; composite
indicators; reserve life
\end{abstract}

% --------------------------------------------------------- reader's note --
\noindent\fbox{\begin{minipage}{0.96\textwidth}\small
\textbf{Who this paper is for, and in what order.} The verdicts, in thirty
seconds: the title, this note, and \S\ref{sec:applied}. The argument, in three
minutes: \S1, \S6, \S10.1, \S11. The machinery, in an hour: \S2 (the ledger),
\S3 (three certificates), \S4 (what the closed ledger guarantees), \S5
(readouts and the componentwise deficit), \S6--\S7 (depletion time and its
passage semantics), \S8 (domain templates at registered status), \S9 (the
interface with institutional delay dynamics). Readers who want only the applied
conclusion can take \S6.5 and the classification table; the theorems are what
license those readings, and nothing else in this paper is required to accept
them.

\smallskip
\textbf{What this paper is not.} It is not a forecast, not an index, and not a
recommendation about any particular aquifer, fishery, or mine. It is an
accounting layer whose purpose is to make sure that when a statement about
depletion is made, the statement that is true is the statement that is written
down.
\end{minipage}}

% ------------------------------------------------- opening of \S1.1 (new) --
\subsection{Two ways the bookkeeping goes wrong}\label{sec:twoways}
A water board announces that a district has a decade or two of groundwater at
current rates. A geological survey lists phosphate reserves of roughly
$1{,}000{,}000$\,kt, and someone divides by production and announces three
centuries. A fisheries assessment divides a log biomass margin by a fishing
mortality and calls the quotient a time scale. All three numbers are computed
correctly. All three answer different questions, and only the last is even
described in the language of its own construction.

Sustainability accounting fails in two characteristic ways. The first is
\emph{compensatory aggregation}: heterogeneous physical stocks and service
flows are summarised by scalar indices whose cross-component trades are never
declared as mathematics, so a severe deficit in one component can coexist with
a positive aggregate. The composite-indicator and
weak-versus-strong-sustainability literatures document the failure and its
noncompensatory remedies (Munda and Nardo, 2009; Ekins et al., 2003; Neumayer,
2013). The second is \emph{classification drift}: quantities carrying the units
of time circulate as if they were one quantity, ``time to depletion'', though
they answer different questions under different assumptions.

We do not dispute the arithmetic of any of these indicators. We dispute the
shared label. The two theses are not equivalent, and only the second is
defensible for a ledger of this kind.

% --------------------------------------------------- gaps named early, \S1 --
\paragraph{Four limits, registered rather than conceded later.}
(i) The groundwater two-pool model is an open gap; the applied object we admit
is its one-pool affine approximation (\S8.2). (ii) Spawning biomass is not an
abiotic support pool, and we do not treat it as one (\S6.5.4). (iii) Every
first-passage result is conditioned on a declared surrogate and carries
explicit non-claims (\S7.7). (iv) The numerical exhibits are worked instances of
the constructions, not constitutive claims for the named domains. A reader
looking for guidance about a specific resource will not find it here; a reader
looking to know which of their claims their data can carry will.

% --------------------------------------------------- section prefaces (\S) --
% One italic sentence per section; adjust \ref labels to your own.
\newcommand{\preface}[1]{\emph{#1}\par\smallskip}

\preface{We fix the objects first: what a compartment is, what a primitive flux
is, and why a conversion may appear only as an explicit coefficient. Everything
that follows is a property of this bookkeeping and nothing else.}   % \S2

\preface{Three predicates get conflated in the literature: the ledger balances,
the chemistry closes, the barriers hold. We prove each separately so that a
reader can tell which one a given argument uses.}                   % \S3

\preface{Here the ledger pays for itself. Mass is conserved because of the
incidence pattern, nothing goes negative because every outflow is
donor-limited, and no positive-effort interior rest exists to hide in.}  % \S4

\preface{Services are read off the ledger, not added to it. The deficit is a
vector, and \S5.4 shows exactly where a scalar would have hidden it.}  % \S5

\preface{Three quantities, three questions. We keep them named separately through
\S6.5, where each public indicator is classified at the status its construction
supports.}                                                             % \S6

\preface{A hitting time is a property of a model, so we state the model, the
barrier, and the conditioning, and then say what does not follow (\S7.7).}  % \S7

\preface{Templates for phosphorus, groundwater, and harvest-side economics, each
with its status and its admitted object marked, so a reader can see what is a
construction and what is a claim about the world.}                    % \S8

\preface{One shared object, two systems, and a non-reduction result that says
why the connection is an interface contract rather than a limit.}     % \S9

\preface{Not caveats: prohibitions on a class of argument. Five double-counting
rules, a theorem against compensatory weighting, and a discipline that keeps
negative and boundary content first-class.}                          % \S10

% --------------------------------------- reusable verdict sentences (\S6.5) --
\newcommand{\verdictGThreeP}{A trend index on a gravity anomaly is a statistical
statement about a record, not a stock ratio.}
\newcommand{\verdictPhosphate}{Reserves divided by production is arithmetic on an
economic classification, not a forecast; the classification refills itself as it
is drawn down.}
\newcommand{\verdictFisheries}{A log margin divided by a removal rate is a
pressure scale, not a time to any event.}
% Use \verdictGThreeP\ in \S1, in the \S6.5 heading text, and again in \S11.

% ------------------------------------------------------ \S11 closing lines --
Nothing in this paper argues that the numbers are wrong. It argues that a reader
should be able to tell, from the statement itself, which question was answered.
That is a bookkeeping property, and it is cheap to install.

%% ------------------------------------------------------------------- demo --
%% \documentclass{article}
%% \usepackage{amsmath,amssymb}
%% \newlabel{sec:applied}{{6.5}{0}}
%% \begin{document}
%% %% snippets.tex --- LaTeX-ready blocks for the humanization pass.
%% Self-contained: paste the pieces you want. Uses only standard article +
%% amsmath; the boxes use \fbox/minipage so no extra packages are required.
%% Check by appending the "demo" document at the end and running pdflatex.

% ---------------------------------------------------------------- abstract --
\begin{abstract}
Someone publishes a number in years --- an aquifer with a decade or two left, a
phosphate reserve of three centuries --- and a reader hears a date. Three
different quantities travel under that single label: a ratio of an economic
reserve classification to current production, a trend fitted to a satellite
gravity anomaly, and a log biomass margin divided by a fishing mortality. Each
answers a real question; none is a time at which anything physically runs out.
A weighted aggregate compounds the confusion by hiding a component deficit
behind a positive scalar.

We build the accounting layer that keeps the three apart. Compartments carry a
material identity, a boundary, and a unit, so biomass, money, and biodiversity
are never summed into one scalar; conservation follows from the incidence
structure of the flux network, positivity from donor limitation, and services
are readouts rather than conserved mass. We separate three predicates the
literature conflates --- accounting consistency, stoichiometric conservation,
and barrier admissibility --- and prove their relations, so that each claim about
a material system carries the predicate it actually establishes. Depletion time
splits into gross turnover intensity, a frozen-rate ratio, and a
scenario-conditioned hitting time, with uniform-drift bounds between them. No
nonnegative weighting of component balances certifies that every component
clears its floor: certification needs the vector. We classify three public
applications at their exact status and state the surrogate on which each
first-passage result is conditioned. We name the gaps that remain open, and we
read weak and strong sustainability as two regimes of one system, distinguished
by whether the material cycle closes at the rate of use.

\medskip
\noindent\textbf{Keywords:} material flow accounting; stock--flow ledger;
depletion indicators; first-passage time; conservation laws; composite
indicators; reserve life
\end{abstract}

% --------------------------------------------------------- reader's note --
\noindent\fbox{\begin{minipage}{0.96\textwidth}\small
\textbf{Who this paper is for, and in what order.} The verdicts, in thirty
seconds: the title, this note, and \S\ref{sec:applied}. The argument, in three
minutes: \S1, \S6, \S10.1, \S11. The machinery, in an hour: \S2 (the ledger),
\S3 (three certificates), \S4 (what the closed ledger guarantees), \S5
(readouts and the componentwise deficit), \S6--\S7 (depletion time and its
passage semantics), \S8 (domain templates at registered status), \S9 (the
interface with institutional delay dynamics). Readers who want only the applied
conclusion can take \S6.5 and the classification table; the theorems are what
license those readings, and nothing else in this paper is required to accept
them.

\smallskip
\textbf{What this paper is not.} It is not a forecast, not an index, and not a
recommendation about any particular aquifer, fishery, or mine. It is an
accounting layer whose purpose is to make sure that when a statement about
depletion is made, the statement that is true is the statement that is written
down.
\end{minipage}}

% ------------------------------------------------- opening of \S1.1 (new) --
\subsection{Two ways the bookkeeping goes wrong}\label{sec:twoways}
A water board announces that a district has a decade or two of groundwater at
current rates. A geological survey lists phosphate reserves of roughly
$1{,}000{,}000$\,kt, and someone divides by production and announces three
centuries. A fisheries assessment divides a log biomass margin by a fishing
mortality and calls the quotient a time scale. All three numbers are computed
correctly. All three answer different questions, and only the last is even
described in the language of its own construction.

Sustainability accounting fails in two characteristic ways. The first is
\emph{compensatory aggregation}: heterogeneous physical stocks and service
flows are summarised by scalar indices whose cross-component trades are never
declared as mathematics, so a severe deficit in one component can coexist with
a positive aggregate. The composite-indicator and
weak-versus-strong-sustainability literatures document the failure and its
noncompensatory remedies (Munda and Nardo, 2009; Ekins et al., 2003; Neumayer,
2013). The second is \emph{classification drift}: quantities carrying the units
of time circulate as if they were one quantity, ``time to depletion'', though
they answer different questions under different assumptions.

We do not dispute the arithmetic of any of these indicators. We dispute the
shared label. The two theses are not equivalent, and only the second is
defensible for a ledger of this kind.

% --------------------------------------------------- gaps named early, \S1 --
\paragraph{Four limits, registered rather than conceded later.}
(i) The groundwater two-pool model is an open gap; the applied object we admit
is its one-pool affine approximation (\S8.2). (ii) Spawning biomass is not an
abiotic support pool, and we do not treat it as one (\S6.5.4). (iii) Every
first-passage result is conditioned on a declared surrogate and carries
explicit non-claims (\S7.7). (iv) The numerical exhibits are worked instances of
the constructions, not constitutive claims for the named domains. A reader
looking for guidance about a specific resource will not find it here; a reader
looking to know which of their claims their data can carry will.

% --------------------------------------------------- section prefaces (\S) --
% One italic sentence per section; adjust \ref labels to your own.
\newcommand{\preface}[1]{\emph{#1}\par\smallskip}

\preface{We fix the objects first: what a compartment is, what a primitive flux
is, and why a conversion may appear only as an explicit coefficient. Everything
that follows is a property of this bookkeeping and nothing else.}   % \S2

\preface{Three predicates get conflated in the literature: the ledger balances,
the chemistry closes, the barriers hold. We prove each separately so that a
reader can tell which one a given argument uses.}                   % \S3

\preface{Here the ledger pays for itself. Mass is conserved because of the
incidence pattern, nothing goes negative because every outflow is
donor-limited, and no positive-effort interior rest exists to hide in.}  % \S4

\preface{Services are read off the ledger, not added to it. The deficit is a
vector, and \S5.4 shows exactly where a scalar would have hidden it.}  % \S5

\preface{Three quantities, three questions. We keep them named separately through
\S6.5, where each public indicator is classified at the status its construction
supports.}                                                             % \S6

\preface{A hitting time is a property of a model, so we state the model, the
barrier, and the conditioning, and then say what does not follow (\S7.7).}  % \S7

\preface{Templates for phosphorus, groundwater, and harvest-side economics, each
with its status and its admitted object marked, so a reader can see what is a
construction and what is a claim about the world.}                    % \S8

\preface{One shared object, two systems, and a non-reduction result that says
why the connection is an interface contract rather than a limit.}     % \S9

\preface{Not caveats: prohibitions on a class of argument. Five double-counting
rules, a theorem against compensatory weighting, and a discipline that keeps
negative and boundary content first-class.}                          % \S10

% --------------------------------------- reusable verdict sentences (\S6.5) --
\newcommand{\verdictGThreeP}{A trend index on a gravity anomaly is a statistical
statement about a record, not a stock ratio.}
\newcommand{\verdictPhosphate}{Reserves divided by production is arithmetic on an
economic classification, not a forecast; the classification refills itself as it
is drawn down.}
\newcommand{\verdictFisheries}{A log margin divided by a removal rate is a
pressure scale, not a time to any event.}
% Use \verdictGThreeP\ in \S1, in the \S6.5 heading text, and again in \S11.

% ------------------------------------------------------ \S11 closing lines --
Nothing in this paper argues that the numbers are wrong. It argues that a reader
should be able to tell, from the statement itself, which question was answered.
That is a bookkeeping property, and it is cheap to install.

%% ------------------------------------------------------------------- demo --
%% \documentclass{article}
%% \usepackage{amsmath,amssymb}
%% \newlabel{sec:applied}{{6.5}{0}}
%% \begin{document}
%% \input{snippets.tex}
%% \end{document}

%% \end{document}

%% \end{document}

%% \end{document}
